% Part: incompleteness
% Chapter: incompleteness-provability
% Section: provability-conditions

\documentclass[../../../include/open-logic-section]{subfiles}

\begin{document}

\olfileid{inc}{inp}{prc}

\olsection{شرایط \usetoken{S}{derivability} برای~$\Th{PA}$}

حساب پئانو، یا~$\Th{PA}$، نظریه‌ای است که~$\Th{Q}$ را با اصول استقرا
برای همهٔ !!{formula}s گسترش می‌دهد. به بیان دیگر، به~$\Th{Q}$ اصولی
به صورت
\[
(!A(0) \land \lforall[x][(!A(x) \lif !A(x'))]) \lif \lforall[x][!A(x)]
\]
برای هر !!{formula}~$!A$ افزوده می‌شوند. توجه کنید که این
در واقع یک \emph{شِما} است، یعنی بی‌نهایت اصل موضوع (و معلوم می‌شود که
$\Th{PA}$ به‌طور متناهی اصل‌موضوعی‌پذیر {\em نیست}). اما چون می‌توان
به‌طور مؤثر تعیین کرد که آیا رشته‌ای از نمادها نمونه‌ای از یک اصل استقرا
است یا نه، مجموعهٔ اصول موضوع~$\Th{PA}$ محاسبه‌پذیر است. $\Th{PA}$
نظریه‌ای بسیار نیرومندتر از~$\Th{Q}$ است. برای نمونه، با به‌کارگیری
استقرا به شیوهٔ معمول، به‌آسانی می‌توان اثبات کرد که جمع و ضرب جابه‌جایی
هستند. در واقع، بیشتر استدلال‌های متناهیِ نظریهٔ اعداد و ترکیبیاتی را
می‌توان در~$\Th{PA}$ انجام داد.

از آنجا که~$\Th{PA}$ به‌طور محاسبه‌پذیر اصل‌موضوعی‌سازی شده است، محمول
!!{derivability} یعنی~$\Prf[\Th{PA}](x,y)$ محاسبه‌پذیر است و بنابراین
در~$\Th{Q}$ (و در نتیجه در~$\Th{PA}$) بازنمایی می‌شود. مانند پیش،
$\OPrf[\Th{PA}](x,y)$ را برای فرمول بازنمایندهٔ این رابطه به کار
می‌بریم. بگذارید~$\OProv[\Th{PA}](y)$ فرمول
$\lexists[x][\OPrf[\Th{PA}](x,y)]$ باشد که، به‌طور شهودی، می‌گوید
«$y$ از اصول موضوع~$\Th{PA}$ !!{derivable} است.» علت اینکه به اندکی
بیش از اصول موضوع~$\Th{Q}$ نیاز داریم این است که باید بدانیم نظریه‌ای
که به کار می‌بریم آن‌قدر نیرومند هست که چند واقعیت بنیادی را !!{derive}
کند؛ واقعیت‌هایی دربارهٔ این محمول !!{derivability}. در واقع، به
واقعیت‌های زیر نیاز داریم:
\begin{enumerate}
\item[P1.] اگر~$\Th{PA} \Proves !A$، آنگاه~$\Th{PA} \Proves
  \OProv[\Th{PA}](\gn{!A})$.
\item[P2.] برای همهٔ !!{formula}s چون~$!A$ و~$!B$،
  \[
  \Th{PA} \Proves \OProv[\Th{PA}](\gn{!A \lif !B}) \lif
  (\OProv[\Th{PA}](\gn{!A}) \lif \OProv[\Th{PA}](\gn{!B})).
  \]
\item[P3.] برای هر !!{formula}~$!A$،
  \[
  \Th{PA} \Proves \OProv[\Th{PA}](\gn{!A})
  \lif \OProv[\Th{PA}](\gn{\OProv[\Th{PA}](\gn{!A})}).
  \]
\end{enumerate}
تنها راه احراز این سه ویژگی آن است که !!{formula}
$\OProv[\Th{PA}](y)$ را به‌دقت توصیف کنیم و با استفاده از اصول
$\Th{PA}$، !!{derivation}s صوری مربوط را شرح دهیم. شرط‌های (۱) و~(۲)
آسان‌اند؛ در واقع شرط~(۳) است که کار می‌طلبد. (بیندیشید که این کار چه
نوع کاری است~\dots) پرداختن به جزئیات ملال‌آور و بی‌فایده خواهد بود،
پس اینجا از شما می‌خواهیم با اعتماد بپذیرید که~$\Th{PA}$ سه ویژگی
فهرست‌شده در بالا را دارد. انتخابی معقول برای~$\OProv[\Th{PA}](y)$
ویژگی زیر را نیز برآورده خواهد کرد:
\begin{enumerate}
\item[P4.] اگر~$\Th{PA} \Proves \OProv[\Th{PA}](\gn{!A})$، آنگاه
  $\Th{PA} \Proves !A$.
\end{enumerate}
اما به این واقعیت نیاز نخواهیم داشت.

\begin{digress}
گفتنی است که گودل نیز از همان جهتی تن‌آسا بود که ما اکنون هستیم. او در
پایان مقالهٔ سال ۱۹۳۱ طرح اثبات قضیهٔ دوم ناتمامیت را ارائه می‌کند و
وعده می‌دهد جزئیات را در مقاله‌ای بعدی بیاورد. هرگز به انجام آن نرسید؛
زیرا همهٔ کسانی که استدلال را می‌فهمیدند باور داشتند که می‌توان آن را
کامل کرد (او نیازی به پر کردن جزئیات نداشت.)
\end{digress}

\end{document}
